\documentclass[conference]{IEEEtran}

\usepackage{cite}
\usepackage{amsmath,amssymb,amsfonts}
\usepackage{graphicx}
\usepackage{booktabs}
\usepackage{array}
\usepackage{url}
\usepackage{xcolor}

\title{Notes on Tuning the Hyperopt Objective Function in Practice}

\author{
\IEEEauthorblockN{Luis Mendez}
}

\begin{document}
\maketitle

\section*{Overview}
The companion note (\textit{Other SMBO Models}) covers what \texttt{rand.suggest}, \texttt{anneal.suggest}, and \texttt{tpe.suggest} do algorithmically. This note is about a narrower, more practical problem that shows up once those algorithms are pointed at a real objective function: \texttt{02-Search-algorithms.ipynb} tunes an \texttt{XGBClassifier} on the breast cancer dataset with all three algorithms, and in every run the trained model reaches a train ROC-AUC of exactly $1.0$, while a 50-trial TPE search took 3h22m to finish. Neither number is intrinsic to the search algorithm; both trace back to how the objective function was written. Fixing the objective, not the search algorithm, is what moved the needle.

\section*{Diagnosis 1: CPU Oversubscription}
The objective function scored each candidate with
\begin{verbatim}
cross_val_score(gbm, X_train, y_train,
    scoring='accuracy', cv=3, n_jobs=4)
\end{verbatim}
which forks 4 worker processes, one per fold-batch. Each worker holds an \texttt{XGBClassifier} built with its \emph{default} thread count, i.e.\ ``use every core available.'' Four processes each trying to claim every core means the OS is constantly context-switching between more threads than there are physical cores to run them on --- classic oversubscription. This is consistent with the wall-clock numbers observed: random search (cheaper, shallower trees on average) finished in 52 minutes, annealing was still running after 46 minutes at trial 28/50, and TPE took 3h22m for the same 50 trials. The fix is to pin the model itself to a single thread and let \texttt{cross\_val\_score} own all the parallelism:
\begin{verbatim}
params_dict = {
    ...
    'n_jobs': 1,
    'tree_method': 'hist',
}
\end{verbatim}
\texttt{tree\_method='hist'} (histogram-binned split finding) is also substantially faster per tree than XGBoost's exact-greedy default, independent of the threading fix. Together, these make each trial cheaper, which is what actually buys a bigger effective search --- not raising \texttt{max\_evals} against the original, oversubscribed objective.

\section*{Diagnosis 2: An Unregularized Corner of the Search Space}
A train ROC-AUC of $1.0$ on every single trial, across three different search algorithms and search-space regions, is a sign the model has enough capacity to perfectly separate every training point regardless of which hyperparameters are drawn. The original space,
\begin{equation}
    \{n_{\text{est}}, \text{depth}, \eta, \text{booster}, \gamma, \text{subsample}, \text{colsample}_{*}, \lambda\},
\end{equation}
constrains \emph{how much} data and how many features each tree sees, but not \emph{how small a leaf is allowed to be} or \emph{how large a raw split-weight is allowed to grow}. Two parameters close that gap:
\begin{itemize}
    \item \textbf{\texttt{min\_child\_weight}} (searched over $[1, 20]$): the minimum sum of instance Hessian weight required in a leaf before XGBoost will keep splitting. Low values let the tree carve out leaves around single outlying training points; raising the floor forces leaves to aggregate enough evidence to be trustworthy.
    \item \textbf{\texttt{reg\_alpha}} (searched over $[0, 5]$): the $L_1$ penalty on leaf weights, complementing the $L_2$ penalty already present as \texttt{reg\_lambda}. $L_1$ additionally pushes small leaf weights to exactly zero, pruning marginal splits that $L_2$ alone would only shrink.
\end{itemize}
Neither guarantees train AUC drops below $1.0$ --- with $n_{\text{estimators}}$ up to 2500 the model can still memorize --- but the search is now able to \emph{find} more conservative regions of the space if those regions generalize better, which it could not do before because those regions did not exist in the space at all.

\section*{Diagnosis 3: Noisy 3-Fold Scoring}
With \texttt{cv=3} and no fixed fold assignment, two evaluations of the same hyperparameter set can land on different folds and disagree purely from sampling noise, on a dataset with only 398 training rows. Since \texttt{anneal.suggest} and \texttt{tpe.suggest} both use the trial history to decide where to sample next, they are partly reacting to that noise rather than to genuine differences between hyperparameter sets. Two changes address this:
\begin{verbatim}
cv = StratifiedKFold(
  n_splits=5, shuffle=True,
  random_state=0)
\end{verbatim}
fixed once, outside the objective, and reused for every trial. Five folds instead of three shrinks the variance of the mean CV score; fixing the split (rather than letting \texttt{cross\_val\_score} re-derive a fold assignment implicitly) means every trial is scored on exactly the same train/validation partitions, so differences in the reported score are attributable to the hyperparameters and not to which rows happened to land in which fold.

\section*{A Free Variance-Reduction Step: Ensembling the Three Searches}
\texttt{rand.suggest}, \texttt{anneal.suggest}, and \texttt{tpe.suggest} were run independently and each produces its own locally-optimal configuration; they do not agree exactly, since each is a different (partly stochastic) path through the same space. Rather than picking one, averaging the three fitted models' predicted probabilities,
\begin{equation}
    \hat{p}(\mathbf{x}) = \frac{1}{3}\Big(\hat{p}_{\text{rand}}(\mathbf{x}) + \hat{p}_{\text{anneal}}(\mathbf{x}) + \hat{p}_{\text{tpe}}(\mathbf{x})\Big),
\end{equation}
is a zero-cost variance-reduction step at prediction time: no extra trials, no extra tuning, just three models whose errors are not perfectly correlated (each converged to a different region of hyperparameter space) getting averaged together. This is the same intuition behind bagging, applied post-hoc to independently-tuned models instead of independently-resampled datasets.

\section*{Practical Guidance}
\begin{itemize}
    \item Before increasing \texttt{max\_evals} to search harder, check whether the objective function itself is leaving performance (in this case, wall-clock time) on the table --- nested parallelism between \texttt{cross\_val\_score} and the model's own thread count is an easy one to miss.
    \item A train metric stuck at a perfect score across every trial is a search-space smell, not just a modeling one: it means the space never offered the search algorithm a sufficiently regularized configuration to try.
    \item Fix the cross-validation folds once, outside the objective, whenever a search algorithm uses trial history to decide where to sample next --- otherwise part of what it is ``learning'' is fold-assignment noise.
    \item When multiple search algorithms are run over the same space anyway (e.g.\ to compare them, as here), averaging their resulting models is a free way to use all of that compute instead of discarding two of the three results.
\end{itemize}

\end{document}
